\section*{2. 关于水平的回顾}
\noindent(参见~[2], 9.7 和 10.1).

\subsubsection*{2.0.1.}
设 $X$ 在代数闭域 $k$ 上固有且光滑. 在 (1.2.3) 中已经看到, 对每个
\begin{equation}\tag{2.0.2}
T\quad\text{在 $k$ 上固有、光滑且连通, 维数为 $d$},
\end{equation}
和每个
\begin{equation}\tag{2.0.3}
\text{$T\times X$ 上余维为 $b$ 的代数圈 $z$},
\end{equation}
上同调类
\begin{equation}\tag{2.0.4}
\mathrm{cl}(z)\in H^{2b}(T\times X,\mathbb Q_\ell(b))
\end{equation}
给出映射
\begin{equation}\tag{2.0.5}
\mathrm{cl}(z):H^q(T,\mathbb Q_\ell(i))\longrightarrow
H^{q+2b-2d}(X,\mathbb Q_\ell(i+b-d)).
\end{equation}
$H^q(X,\mathbb Q_\ell(i))$ 上的\underline{余水平滤过}(coniveau filtration) $(N^j)$ 定义为
\begin{equation}\tag{2.0.6}
N^jH^q(X,\mathbb Q_\ell(i))
=\sum_{(T,z)}\operatorname{Im}\!\left(
\mathrm{cl}(z):H^{q-2j}(T,\mathbb Q_\ell(i-j))
\longrightarrow H^q(X,\mathbb Q_\ell(i))\right),
\end{equation}
其中求和遍历所有对
\begin{equation}\tag{2.0.7}
(T,z),\qquad
\begin{array}{l}
T/k\text{ 固有、光滑且连通, 维数为 }d,\\
z\text{ 为 }T\times X\text{ 上余维为 }j+d\text{ 的代数圈}.
\end{array}
\end{equation}
有时也采用``按水平''的递增滤过 $F^j$, 其定义为
\begin{equation}\tag{2.0.8}
F^jH^q(X,\mathbb Q_\ell(i))
=N^{\left[\frac{q-j}{2}\right]}H^q(X,\mathbb Q_\ell(i)).
\end{equation}
用 $(\mathbb P^1\times T,\{\mathrm{pt}\}\times z)$ 代替 $(T,z)$, 可见
\begin{equation}\tag{2.0.9}
N^jH^q(X,\mathbb Q_\ell(i))
\subset N^{j-1}H^q(X,\mathbb Q_\ell(i)).
\end{equation}
显然有:

\subsubsection*{命题 2.1.}
设 $X$ 在域 $k$ 上固有、光滑且几何连通, $\bar k$ 为 $k$ 的代数闭包, 置 $X_{\bar k}=X\otimes_k\bar k$. 则 $\operatorname{Gal}(\bar k/k)$ 在 $H^q(X_{\bar k},\mathbb Q_\ell(i))$ 上的作用保持余水平滤过.

标准论证 (参见~[2]) 可证明下列断言.

\subsubsection*{命题 2.2.}
水平滤过在代数闭域的扩张 $k'/k$ 下不变.

\subsubsection*{命题 2.3.}
设 $S$ 为约化不可约概形, 泛点为 $\eta$, 并设
\begin{equation*}
\pi:X\longrightarrow S
\end{equation*}
为具有几何连通纤维的固有光滑态射. 若对于给定整数 $j$ 和 $q$, 在几何泛纤维 $X_{\bar\eta}$ 上有
\begin{equation}\tag{2.3.1}
H^q(X_{\bar\eta},\mathbb Q_\ell)
=N^jH^q(X_{\bar\eta},\mathbb Q_\ell),
\end{equation}
则存在非空开子集 $U\subset S$, 使得对于每个点 $u\in U$, 在几何纤维 $X_{\bar u}$ 上有
\begin{equation}\tag{2.3.2}
H^q(X_{\bar u},\mathbb Q_\ell)
=N^jH^q(X_{\bar u},\mathbb Q_\ell).
\end{equation}

